\section{Delayed information needs are absent from AITZ and unlabeled in GUI-Odyssey}
\label{sec:needs}
A step $t$ has a \emph{delayed information need} when the text the agent types at $t$ appears verbatim in the OCR text of an earlier screen $i<t$ (a \emph{source screen}) and does not appear on screen $t$. The \emph{gap} is $t-i$ for the latest such $i$; earliest-source scoring (\S\ref{sec:pitfalls}) instead labels the earliest. The second clause of the definition matters: when the text is on the current screen the agent needs no memory, and \S\ref{sec:pitfalls} shows that keeping such steps corrupts every comparison.

We mine candidates with a string rule and verify them with OCR. A candidate is any TYPE action whose argument has three or more characters and does not occur in the task description; we normalize strings by lower-casing and stripping whitespace and punctuation. We run EasyOCR on every screenshot of 300 sampled GUI-Odyssey trajectories and test whether the argument is a substring of an earlier screen's OCR text. Of 431 candidates, 127 pass (304 fail: 269 where no earlier screen's OCR text contains the string, most of them OCR misreads on inspection, and 35 where OCR output is missing). Removing 36 whose text is also on the current screen leaves 91 verified cases; removing 6 whose text appears in the task's step-by-step instruction field leaves the 85-case short set. The count is a lower bound, and the rule covers only verbatim re-typing. Table~\ref{tab:prev} gives prevalence from the 91 verified cases, 89 of which fall in the 273 sampled trajectories with complete OCR coverage.

\begin{table}[t]
\centering\footnotesize
\caption{Delayed information needs are absent from AITZ and present in a third of GUI-Odyssey trajectories.}
\label{tab:prev}
\begin{tabular}{lrrrl}
\toprule
 & trajectories & steps & mean length & verified delayed needs \\
\midrule
AITZ~\citep{aitz} & 2,504 & 18,643 & 7.4 & 0 \\
GUI-Odyssey~\citep{guiodyssey} (full) & 7,735 & 118,791 & 15.4 & \\
\quad sampled, complete OCR & 273 & 4,338 & 15.9 & 89 steps (2.05\%) in 85 of the 273 trajectories (31.1\%) \\
\bottomrule
\end{tabular}
\end{table}

The need is common, far, and inconspicuous. The gap has median 5 and maximum 27 steps, and 56\% of gaps are 5 or more. On its source screen the needed text sits at a median rank of 13 among 33 OCR elements and is never first. A one-line summary that keeps the first 3--5 elements in reading order therefore retains the answer in 7--20\% of cases (the top-3 and top-5 rates) before any reader sees it.

We evaluate on three case sets, each filtered as above, which we call \emph{clean}. The \emph{short set} holds the 85 cases above (the need's step index $t$ has median 11). The \emph{long set} holds 53 cases with $t\ge16$ from a second sample of 300 trajectories of 20 or more steps (by coincidence also 127 before the leakage filter, all 74 removals being leakage). The instruction-field check flags one further long-set case, a substring artifact, which we keep. The \emph{held-out set} holds 50 cases from a third sample of 200 trajectories, mined with the final filter and evaluated once with every setting fixed (\S\ref{sec:length}).
